\section{The method}\label{sec:method}

Rite splits into an engine and per-paper workdirs (Figure~\ref{fig:sys}).
The engine is read-only at writing time; this is a trust boundary as much as a convenience, since a tool that edits its own rules while writing paper seven will eventually rot them.

The pipeline is disciplined by two single sources of truth.
The search goal and year range live only in the workdir README; the coding vocabulary (topic flags and technology flags, regular-expression families) lives only in the workdir \texttt{flags.py}.
Retargeting the whole pipeline to a new paper therefore means editing just two files.

The structure division pins down the shape of the paper itself.
Figure~\ref{fig:bnf} states that shape as a definite clause grammar, in the Prolog idiom where each head argument is a \texttt{key=Variable} pair: the keys document which role each part plays, and any part of a parsed paper can be pulled out by name.
The keywords earn their keep where one nonterminal is called twice in a single rule (in \texttt{related}, backward and forward snowballing are both just \texttt{cites}); there, position alone would be a silent trap.

\begin{figure}[t]
\begin{lstlisting}
paper(front=F, body=B, back=K) -->
    front(F), body(B), back(K).

front(title=T, author=A, abstract=X) -->
    title(T), author(A), abstract(X).

body(intro=I, related=R, method=M,
     results=E, discussion=D,
     conclusion=C) -->
    intro(I), related(R), method(M),
    results(E), discussion(D),
    conclusion(C).

intro(quote=Q, questions=Qs,
      contributions=Cs, faq=A,
      install=In, use=U, collab=V) -->
    quote(Q), rqs(Qs), contribs(Cs),
    faq(A), install(In), use(U),
    collab(V).

related(before=B, since=S) -->
    cites(B),  % backward snowball
    cites(S).  % forward snowball

discussion(observations=O, threats=T,
           future=W) -->
    observations(O), threats(T),
    future(W).

back(refs=Rs) --> refs(Rs).
\end{lstlisting}
\caption{The shape of a standard paper, as the structure division states it: a definite clause grammar whose head arguments name their parts.}
\label{fig:bnf}
\end{figure}

The pipeline then runs from the workdir: search OpenAlex for the goal; keep the papers above the knee of the sorted-citation curve; backward-snowball their references for the shared classics; fetch open-access PDFs; code abstracts, then code full texts.
Full-text coding counts a flag only above a frequency threshold (0.5 matches per 1000 words), a crude form of length-normalized term weighting~\cite{salton88}; the cutoff value is ours and Section~\ref{sec:disc} owes it a sensitivity analysis.

The loop's most opinionated gate is the abstract-consistency rule (step 13): run the paper's own coding grid on its own title plus abstract, and require every flag the body earns at threshold.
The justification is Section~\ref{sec:exhibit}: in the wild, abstracts fail this test about half the time, and everyone who judges a paper cheaply (bidding PC members, systematic-review coders, ranking algorithms, skimming practitioners) reads only the abstract.
